\subsection{Simple Subalgebras of Simple Algebras}

THEOREM 5. --- Let A be a central simple K-algebra, and let B be a semisimple subalgebra of A of finite degree over K.

\begin{enumerate}
\def\labelenumi{\alph{enumi})}
\item
  The commutant B' of B in A is a simple subalgebra, and B is the commutant of B' in A. Moreover, the algebra \(B \cap B'\) is a semisimple commutative algebra of finite degree over K; it is the common center of B and B'.
\item
  Suppose that B is simple. Then \(B'\) is simple, and we have the equalities
\end{enumerate}

\[
[ \mathrm{A}: \mathrm{B} ^ {\prime} ] _ {s} = [ \mathrm{B}: \mathrm{K} ], \quad [ \mathrm{A}: \mathrm{B} ] _ {s} = [ \mathrm{B} ^ {\prime}: \mathrm{K} ], \quad [ \mathrm{A}: \mathrm{K} ] = [ \mathrm{B}: \mathrm{K} ] [ \mathrm{B} ^ {\prime}: \mathrm{K} ].
\]

(See VIII, p.~124, Definition 2 for the definition of the left degree \([A : B]_s\) .)

The K-algebra \(A^{\circ}\) is central simple, and the K-algebra B is semisimple and of finite degree. By Corollary 1 of VIII, p.~221, the algebra \(C = B \otimes_{K} A^{\circ}\) is semisimple. Let M be the C-module with the same additive group as A and with external law given by the formula \((b \otimes a)a' = ba'a\) for \(a, a'\) in A and b in B. Let u be an element of \(\operatorname{End}_{\mathbf{Z}}(\mathrm{A})\) . Then u belongs to the commutant \(C_{M}'\) of the C-module M if and only if u is right A-linear and left B-linear, in other words, if and only if u belongs to the commutant of \(B_{M}\) in the ring of homotheties of the A-module \(A_{s}\) . We consequently define an isomorphism \(\gamma\) from \(B'\) to \(C_{M}'\) by the relation \(\gamma(b')(x) = b'x\) for \(b'\) in \(B'\) and x in M. Now, the ring C is semisimple, and the C-module M is generated by the element 1 of A. By Proposition 6 of VIII, p.~139, the ring \(C_{M}'\) is semisimple, so the algebra \(B'\) is semisimple.

Let \(\varphi\) be the K-algebra homomorphism from \(A \otimes_K A^\circ\) to \(\mathrm{End}_K(M)\) that sends \(a \otimes a'\) to the K-linear mapping \(x \mapsto axa'\) from M to M. Since the K-algebras A and \(A^\circ\) are central simple, the only two-sided ideals of \(A \otimes_K A^\circ\) are 0 and \(A \otimes_K A^\circ\) (VIII, p.~218, Corollary). We have \(C_M = \varphi(B \otimes A^\circ)\) and \(C_M' = \varphi(B' \otimes K)\) . The homomorphism \(\varphi\) is not zero; it is therefore injective. Since the ring C is semisimple, we have \(C_M'' = C_M\) by Proposition 5 of VIII, p.~139. It follows that the subalgebra \(B \otimes_K A^\circ\) of \(A \otimes_K A^\circ\) is the commutant of the subalgebra \(B' \otimes_K K\) . The commutant of \(B' \otimes_K K\) in \(A \otimes_K K\) is therefore equal to \((B \otimes_K A^\circ) \cap (A \otimes_K K)\) , that is, to \(B \otimes K\) by Proposition 19 of II, §7, No.~9, p.~311. Hence the commutant of \(B'\) in A is equal to B. The algebra \(L = B \cap B'\) is the center of B. Since B is a semisimple algebra of finite degree over K, the algebra L is commutative and semisimple and has finite degree over K. Since B is the commutant of \(B'\) in A, the center of \(B'\) is also equal to \(L = B \cap B'\) (VIII, p.~77). We have proved a).

Now suppose that the algebra B is simple. By Corollary 2 of VIII, p.~222, the ring C is simple. By Proposition 4 of VIII, p.~123 applied to the C-module M, whose commutant is isomorphic to \(B'\) , the ring \(B'\) is simple and M is a \(B'\) -module of finite length. In other words, \(B'\) is a simple subring of the simple ring A, and the left degree \([A : B']_{s}\) is an integer \(m \geqslant 1\) . Viewed as a left \(B'\) -module, A has a finite basis \((a_{1}, \ldots, a_{m})\) . Moreover (loc. cit.), \(\varphi\) induces by restriction an isomorphism from C to \(\mathrm{C}_{\mathrm{M}}'' = \mathrm{End}_{\mathrm{B}'}(\mathrm{A})\) , and the mapping \(c \mapsto (ca_{1}, \ldots, ca_{m})\) from C to \(A^{m}\) is therefore bijective. Consequently, C is a free right A-module of dimension m. Now, we have \(C = B \otimes A^{o}\) , so C is a free right A-module of dimension {[}B : K{]}; it follows that \([A : B']_{s} = m = [B : K]\) . From Proposition 6 of VIII, p.~125, we deduce

\[
[ \mathrm{A}: \mathrm{K} ] = [ \mathrm{A}: \mathrm{B} ^ {\prime} ] _ {s} [ \mathrm{B} ^ {\prime}: \mathrm{K} ] = [ \mathrm{B}: \mathrm{K} ] [ \mathrm{B} ^ {\prime}: \mathrm{K} ];
\]

since we also have

\[
[ \mathrm{A}: \mathrm{K} ] = [ \mathrm{A}: \mathrm{B} ] _ {s} [ \mathrm{B}: \mathrm{K} ]
\]

and \([B:K]\) is finite and nonzero, we conclude that we have the equality \([A:B]_{s}=[B^{\prime}:K]\) (Set Theory, III, §6, No.~3, p.~188, Corollary 3). We have proved b).

\pandocbounded{\includegraphics[keepaspectratio]{images/part002_69e19a74.jpg}}

Let A be a central simple K-algebra of finite degree. There can exist semisimple commutative subalgebras B of A satisfying \([A:K]\neq[B:K][B':K]\) (Exercise 1 of VIII, p.~269).

THEOREM 6. --- Let A be a central simple K-algebra, B a subalgebra of A of finite degree, and B' its commutant in A.

\begin{enumerate}
\def\labelenumi{\alph{enumi})}
\item
  Suppose that B is central simple. Then \(B'\) is central simple, and the K-algebra homomorphism \(\theta : B \otimes_{K} B' \to A\) that sends \(b \otimes b'\) to \(bb'\) is an isomorphism.
\item
  Suppose that B is semisimple, and let \(L = B \cap B'\) . Then \(B'\) is a semisimple algebra. The commutant \(L'\) of L in A is a semisimple ring with center L, and the ring homomorphism \(\psi : B \otimes_{L} B' \to L'\) that sends \(b \otimes b'\) to \(bb'\) is an isomorphism.
\end{enumerate}

Let us prove a). If B is central simple, then \(B'\) is central simple by Theorem 5 of VIII, p.~259. Hence the K-algebra \(B \otimes_{K} B'\) is simple (VIII, p.~222, Corollary 2), and the homomorphism \(\theta : B \otimes_{K} B' \to A\) is injective. Now, by the equality of \([A : B']_{s}\) and \([B : K]\) (Theorem 5), the left \(B'\) -modules \(B \otimes_{K} B'\) and A are free of the same finite dimension; they are therefore \(B'\) -modules of the same finite length. By Corollary 2 of II, §1, No.~10, p.~213, \(\theta\) is bijective.

Let us prove b). By Theorem 5, the algebra L is commutative, of finite degree over K, and semisimple. By loc. cit. applied to L, its commutant \(L'\) in A is a semisimple algebra, and L is the commutant of \(L'\) in A, so L is the center of \(L'\) . Since L is the center of the semisimple rings \(L'\) , B, and \(B'\) , we can identify \(L'\) with a finite product of simple rings \(L_{i}'\) (for \(i \in I\) ), so that we have

\[
\mathrm{L} = \prod_ {i \in \mathrm{I}} \mathrm{L} _ {i}, \qquad \mathrm{B} = \prod_ {i \in \mathrm{I}} \mathrm{B} _ {i}, \qquad \mathrm{B} ^ {\prime} = \prod_ {i \in \mathrm{I}} \mathrm{B} _ {i} ^ {\prime},
\]

where \(L_{i}\) is the center of \(L_{i}^{\prime}\) and where \(B_{i}\) and \(B_{i}^{\prime}\) are subalgebras of \(L_{i}^{\prime}\) with center \(L_{i}\) that are each other's commutants in \(L_{i}^{\prime}\) . We view \(L_{i}^{\prime}\) as a central simple algebra over the commutative field \(L_{i}\) and \(B_{i}\) as a central simple \(L_{i}\) -algebra of finite degree. By assertion a), the canonical mapping \(\psi_{i}: B_{i} \otimes_{L_{i}} B_{i}^{\prime} \to L_{i}^{\prime}\) that sends \(b_{i} \otimes b_{i}^{\prime}\) to \(b_{i}b_{i}^{\prime}\) is a ring isomorphism. Now, we can identify \(B \otimes_{L} B^{\prime}\) with \(\prod_{i \in I}(B_{i} \otimes_{L_{i}} B_{i}^{\prime})\) , so that \(\psi\) is the product of the family of mappings \((\psi_{i})_{i \in I}\) . Therefore, \(\psi\) is a ring isomorphism.

COROLLARY. --- Suppose that the field K is algebraically closed and that A is a simple algebra of finite degree over K. Let B be a simple subalgebra of A, and let B' be the commutant of B in A. Then B' is a simple K-algebra, B is the commutant of B', we have \([A:K]=[B:K][B':K]\) , and the canonical homomorphism from \(B\otimes_{K}B'\) to A is a K-algebra isomorphism.

Since every simple algebra of finite degree over \(K\) is central, the corollary follows from Theorems 5 and 6.

